\documentclass[../../../../main.tex]{subfiles}
\begin{document}

長さ1の線分をつなげてできる右のような平面上の図形 $Q_1,Q_2,Q_3,\cdots$ を考える．
$n=1,2,3,\cdots$ に対し，図形 $Q_n$ の左端の点を $A_n$，右端の点を $B_n$，上端の点を $C_n$ とする．

$Q_1$ は一辺の長さが1の正三角形の周である．
$Q_2$ は図のように，$Q_1$ を3つつなげてできる図形である．

$Q_n$ と同じ図形を3つ用意し，それらを $Q_n(1)$，$Q_n(2)$，$Q_n(3)$ とする．
$i=1,2,3$ に対し，$Q_n(i)$ の左端の点を $A_n(i)$，右端の点を $B_n(i)$，上端の点を $C_n(i)$ としたとき，
$Q_{n+1}$ は，$B_n(1)$ と $A_n(2)$，$C_n(2)$ と $B_n(3)$，$A_n(3)$ と $C_n(1)$ が
それぞれ同一の点になるようにおいできる図形である．

$Q_n$ において，
$A_n$ から線分の上を通り，
一度通った点は二度通らずに $B_n$ まで行く行き方を考える．
この行き方のうち，
途中 $C_n$ を通らない場合の個数を $x_n$ とし，
途中 $C_n$ を通る場合の個数を $y_n$ とする．
容易にわかるように，$x_1=y_1=1$ である．

\begin{figure}[htbp]
  \centering
  \begin{tikzpicture}[scale=0.8]
    % Q1
    \begin{scope}[shift={(0,0)}]
      \draw[thick] (0,0) node[below left] {$A_1$} -- (2,0) node[below right] {$B_1$} -- (1, 1.732) node[above] {$C_1$} -- cycle;
      \node[below=15pt] at (1,0) {図形 $Q_1$};
    \end{scope}

    % Q2
    \begin{scope}[shift={(4,0)}]
      \draw[thick] (0,0) node[below left] {$A_2$} -- (4,0) node[below right] {$B_2$} -- (2, 3.464) node[above] {$C_2$} -- cycle;
      \draw[thick] (2,0) -- (1, 1.732) -- (3, 1.732) -- cycle;
      \node[below=15pt] at (2,0) {図形 $Q_2$};
    \end{scope}
  \end{tikzpicture}
\end{figure}

\begin{enumerate}
  \item $x_2$，$y_2$ を求めよ．
  \item $x_{n+1}$ を $x_n$，$y_n$ を用いて表せ．また，$y_{n+1}$ を $x_n$，$y_n$ を用いて表わせ．
  \item $x_3$，$y_3$ を求めよ．
\end{enumerate}

\end{document}
